% =========================================================
% TP2 -- Automata off-lattice: bandadas de agentes autopropulsados
% Presentacion. Formato segun GuiaPresentaciones.pdf:
%   * diapositivas numeradas (1.2)
%   * poco texto, sin parrafos (1.6)
%   * figuras sin titulo ni epigrafe; los parametros fijos van al costado (1.7)
%   * ejes rotulados en palabras y con unidades entre corchetes (1.8)
%   * notacion cientifica con potencias de 10 (1.9)
%   * sin seccion de bibliografia: cita abreviada en la diapositiva (1.11)
%   * separadores de seccion sin numerar la seccion (1.13)
%   * las animaciones van como cuadro fijo + link explicito (2.4.8)
%
% Las figuras salen de data/figuras_diapositivas/, que es el mismo analisis que
% el del informe pero con tipografia agrandada, para que dentro de la
% diapositiva se lean del tamano del texto de la diapositiva (1.8).
%
%   make figuras-diapositivas
%   tectonic -X compile report/presentacion.tex --outdir report/build
% =========================================================

\documentclass[aspectratio=169,11pt]{beamer}

\usepackage[utf8]{inputenc}
\usepackage[T1]{fontenc}
\usepackage[spanish,es-noquoting,es-noshorthands]{babel}

\usepackage{amsmath}
\usepackage{amssymb}
\usepackage{graphicx}
\usepackage{booktabs}
\usepackage{siunitx}
\sisetup{output-decimal-marker={,}, per-mode=symbol}
\usepackage{tikz}
\usetikzlibrary{arrows.meta,positioning,fit,backgrounds,calc}

% ---------------------------------------------------------
% Tema
% ---------------------------------------------------------

\definecolor{azul}{RGB}{23,66,110}
\definecolor{acento}{RGB}{197,90,17}
\definecolor{gris}{RGB}{110,118,128}
\definecolor{grisclaro}{RGB}{236,238,241}

\usetheme{default}
\usecolortheme{default}
\setbeamertemplate{navigation symbols}{}
\setbeamercolor{structure}{fg=azul}
\setbeamercolor{frametitle}{fg=azul}
\setbeamercolor{title}{fg=azul}
\setbeamercolor{alerted text}{fg=acento}
\setbeamerfont{frametitle}{size=\large,series=\bfseries}
\setbeamerfont{title}{size=\LARGE,series=\bfseries}
\setbeamertemplate{itemize item}{\color{azul}\rule[0.25ex]{0.5ex}{0.5ex}}
\setbeamertemplate{itemize subitem}{\color{gris}\rule[0.25ex]{0.4ex}{0.4ex}}
\setbeamersize{text margin left=0.75cm, text margin right=0.75cm}

% Titulo con una regla fina debajo: separa el titulo del area de datos sin
% ocupar el ancho de un banner de color.
\setbeamertemplate{frametitle}{%
    \vspace*{0.5em}%
    \hspace*{0.75cm}%
    \parbox{0.87\paperwidth}{\raggedright
        \usebeamerfont{frametitle}\usebeamercolor[fg]{frametitle}\insertframetitle\par
        \vspace{-0.2em}\textcolor{azul!30}{\rule{\linewidth}{0.9pt}}}%
}

% Pie: el numero de diapositiva a la derecha (1.2).
\setbeamertemplate{footline}{%
    \hbox{%
    \begin{beamercolorbox}[wd=\paperwidth,ht=2.4ex,dp=1.6ex,leftskip=0.75cm,rightskip=0.75cm]{}%
        {\scriptsize\color{gris}%
         Autómata off-lattice\hfill Grupo 10 -- Comisión S2\hfill\insertframenumber}%
    \end{beamercolorbox}}%
}

% ---------------------------------------------------------
% Atajos
% ---------------------------------------------------------

\newcommand{\va}{\ensuremath{v_a}}
\newcommand{\teq}{\ensuremath{t_{\mathrm{eq}}}}
\newcommand{\Mc}{\ensuremath{M_{\mathrm{c}}}}
% Igual que en el informe: sin depender de la \max que redefine babel.
\newcommand{\Cmax}{\ensuremath{\mathcal{C}_{\mathrm{max}}}}
\newcommand{\cita}[1]{{\scriptsize\color{gris}(#1)}}

% Link a la animacion, debajo del cuadro fijo (2.4.8). El PDF que se entrega no
% lleva animaciones embebidas.
\newcommand{\pendiente}{{\scriptsize\color{gris}\texttt{[link -- pendiente]}}}

% Titulo del bloque de condiciones que acompana a cada figura (1.7).
\newcommand{\setup}[1]{{\color{azul}\bfseries #1}\par\vspace{0.15em}}

% Las cuatro constantes que comparten todas las corridas.
\newcommand{\fijos}{%
    $L = \SI{10}{m}$, contorno periódico\\
    $r_c = \SI{1}{m}$, $v = \SI{0.03}{m/s}$\\
    $\Delta t = \SI{1}{s}$, $M = 20$}

% Figura a la izquierda, condiciones bajo las que se obtuvo a la derecha (1.7).
\newcommand{\figpar}[3][0.66]{%
    \begin{columns}[T,onlytextwidth]
        \begin{column}{#1\textwidth}
            \centering\includegraphics[width=\linewidth]{#2}
        \end{column}\hfill
        \begin{column}{\dimexpr0.97\textwidth-#1\textwidth\relax}
            \raggedright\scriptsize#3
        \end{column}
    \end{columns}}

% Dos figuras y la columna de condiciones: #1 y #2 archivo y rotulo de cada una.
\newcommand{\dosfig}[5]{%
    \begin{columns}[T,onlytextwidth]
        \begin{column}{0.355\textwidth}
            \centering\includegraphics[width=\linewidth]{#1}\\[0.15em]
            {\scriptsize #2}
        \end{column}
        \begin{column}{0.355\textwidth}
            \centering\includegraphics[width=\linewidth]{#3}\\[0.15em]
            {\scriptsize #4}
        \end{column}\hfill
        \begin{column}{0.255\textwidth}
            \raggedright\scriptsize#5
        \end{column}
    \end{columns}}

% Igual que \dosfig pero sin nada debajo de las figuras: los identificadores de
% panel van en el bloque lateral, porque la guia no admite leyendas debajo de la
% figura (1.7). Solo las animaciones llevan texto abajo, que es el link (2.4.8).
\newcommand{\dosfiglimpio}[3]{%
    \begin{columns}[T,onlytextwidth]
        \begin{column}{0.355\textwidth}
            \centering\includegraphics[width=\linewidth]{#1}
        \end{column}
        \begin{column}{0.355\textwidth}
            \centering\includegraphics[width=\linewidth]{#2}
        \end{column}\hfill
        \begin{column}{0.255\textwidth}
            \raggedright\scriptsize#3
        \end{column}
    \end{columns}}

% Identificacion de los dos paneles, arriba del bloque lateral.
\newcommand{\paneles}[2]{%
    {\color{azul}\bfseries Izq.:} #1 \enspace{\color{azul}\bfseries Der.:} #2%
    \par\vspace{0.5em}}

\graphicspath{{../data/figuras_diapositivas/}}

\title{Autómata off-lattice}
\subtitle{Bandadas de agentes autopropulsados}

\begin{document}

% =========================================================
\begin{frame}[plain,noframenumbering]
    \vfill
    \begin{center}
        {\usebeamerfont{title}\usebeamercolor[fg]{title}\inserttitle}\\[0.3em]
        {\large\color{gris}\insertsubtitle}\\[1.0em]
        \textcolor{azul!30}{\rule{0.45\textwidth}{1pt}}\\[1.0em]
        {\small
        \begin{tabular}{ccc}
            Javier Liu & Juan Ignacio Cantarella & Nicolás Rivas \\[0.15em]
            {\footnotesize\color{gris}64332} &
            {\footnotesize\color{gris}64509} &
            {\footnotesize\color{gris}64292}
        \end{tabular}}\\[1.4em]
        {\footnotesize Instituto Tecnológico de Buenos Aires \quad\textbar\quad Simulación de Sistemas}\\[0.3em]
        {\footnotesize\color{gris} Trabajo Práctico 2 \enspace\textbar\enspace Grupo 10 -- Comisión S2 \enspace\textbar\enspace 2026 Q2}
    \end{center}
    \vfill
\end{frame}

% =========================================================
\begin{frame}{Contenido}
    \vspace{1.2em}
    \begin{columns}[T,onlytextwidth]
        \begin{column}{0.5\textwidth}
            \begin{itemize}
                \setlength{\itemsep}{1.0em}
                \item Introducción
                \item Implementación
                \item Simulaciones
            \end{itemize}
        \end{column}
        \begin{column}{0.5\textwidth}
            \begin{itemize}
                \setlength{\itemsep}{1.0em}
                \item Resultados
                \item Conclusiones
            \end{itemize}
        \end{column}
    \end{columns}
\end{frame}

% Separador de seccion: solo el titulo, sin numerar (1.13).
\AtBeginSection[]{%
    \begin{frame}
        \vfill
        \begin{center}
            {\Large\bfseries\color{azul}\insertsectionhead}\\[0.5em]
            \textcolor{azul!30}{\rule{0.28\textwidth}{1pt}}
        \end{center}
        \vfill
    \end{frame}}

% =========================================================
\section{Introducción}
% =========================================================

\begin{frame}{Sistema real}
    \begin{columns}[T,onlytextwidth]
        \begin{column}{0.55\textwidth}
            \vspace{0.8em}
            \begin{itemize}
                \setlength{\itemsep}{1.0em}
                \item Bandadas, cardúmenes, colonias bacterianas: agentes que se
                      \alert{autopropulsan} y se \alert{alinean con sus vecinos}.
                \item Interacción puramente \alert{local}: no hay líder ni campo
                      externo que fije una dirección.
                \item Aun así el conjunto adopta espontáneamente una dirección
                      común.
                \item Un único parámetro de control: el \alert{ruido} $\eta$.
                      Transición continua orden--desorden, fuera del equilibrio.
            \end{itemize}
            \vspace{1.0em}
            \cita{Vicsek et al., Phys. Rev. Lett. 75, 1226, 1995}
        \end{column}\hfill
        \begin{column}{0.41\textwidth}
            \centering
            \includegraphics[width=\linewidth]{a_animaciones/cuadro_vicsek_rho4_eta0.5_t3000.png}
        \end{column}
    \end{columns}
\end{frame}

\begin{frame}{Modelo matemático}
    \vspace{-0.2em}
    {\small $N$ partículas puntuales, rapidez común $v$, dirección $\theta_i$.
    Actualización \alert{sincrónica}: todo en $t+1$ se calcula con el estado en $t$.}
    \vspace{-0.3em}
    \begin{columns}[T,onlytextwidth]
        \begin{column}{0.45\textwidth}
            \begin{equation*}
                \mathbf{r}_i(t+1) = \mathbf{r}_i(t) + \mathbf{v}_i(t)\,\Delta t
            \end{equation*}
        \end{column}
        \begin{column}{0.53\textwidth}
            \begin{equation*}
                \theta_i(t+1) = \Theta_i(t) + \Delta\theta_i,
                \quad
                \Delta\theta_i \sim \mathcal{U}\!\left[-\tfrac{\eta}{2},\tfrac{\eta}{2}\right]
            \end{equation*}
        \end{column}
    \end{columns}
    \vspace{0.2em}
    {\small La \alert{regla de interacción} es lo único que cambia entre los dos
    escenarios. Entorno:
    $\mathcal{N}_i=\{\,j\neq i:\ d(\mathbf{r}_i,\mathbf{r}_j)<r_c\,\}$,
    con $d$ la distancia de mínima imagen.}
    \vspace{0.9em}

    \begin{columns}[T,onlytextwidth]
        \begin{column}{0.48\textwidth}
            {\color{azul}\bfseries Estándar} --- promedia el entorno
            \begin{equation*}
                \Theta_i = \arctan
                \frac{\sin\theta_i + \sum_{j\in\mathcal{N}_i}\sin\theta_j}
                     {\cos\theta_i + \sum_{j\in\mathcal{N}_i}\cos\theta_j}
            \end{equation*}
            \vspace{-0.7em}
            \begin{center}\cita{Vicsek et al., PRL, 1995}\end{center}
        \end{column}\hfill
        \begin{column}{0.48\textwidth}
            {\color{azul}\bfseries Votante} --- copia a uno solo
            \begin{equation*}
                \Theta_i = \theta_k,
                \qquad
                k \sim \mathcal{U}\!\left(\mathcal{N}_i\cup\{i\}\right)
            \end{equation*}
            \vspace{-0.7em}
            \begin{center}\cita{Loscar et al., Phys. Rev. E 104, 034111, 2021}\end{center}
        \end{column}
    \end{columns}
    \vspace{0.7em}
    {\footnotesize $\eta$ en radianes. La propia partícula entra en el promedio y
    en el sorteo, de modo que una partícula aislada conserva su dirección.}
\end{frame}

% =========================================================
\section{Implementación}
% =========================================================

\begin{frame}{Arquitectura del motor}
    \vspace{-0.1em}
    \begin{center}
    \begin{tikzpicture}[
        font=\scriptsize,
        modulo/.style={draw=azul!55,fill=azul!6,rounded corners=2pt,align=center,
                       inner sep=3pt,minimum height=6mm,minimum width=26mm},
        archivo/.style={draw=gris!70,fill=grisclaro,rounded corners=2pt,align=center,
                        inner sep=3pt,minimum height=6mm,minimum width=25mm},
        pyt/.style={draw=acento!65,fill=acento!7,rounded corners=2pt,align=center,
                    inner sep=3pt,minimum height=6mm,minimum width=26mm},
        fl/.style={-{Stealth[length=2mm]},gris,thick}]

        \node[modulo] (arg)  {\texttt{arguments}\\[-2pt]{\tiny parseo y validación}};
        \node[modulo,below=2.2mm of arg] (flock) {\texttt{flock}\\[-2pt]{\tiny estado y regla de actualización}};
        \node[modulo,below=2.2mm of flock] (cim) {\texttt{neighbor\_search}\\[-2pt]{\tiny CIM y fuerza bruta}};
        \node[modulo,below=2.2mm of cim] (obs) {\texttt{observables}\\[-2pt]{\tiny $\va$ y $S$}};
        \node[modulo,below=2.2mm of obs] (io) {\texttt{io}\\[-2pt]{\tiny escritura de texto}};

        \node[archivo,right=15mm of flock] (dyn) {\texttt{dynamic.txt}\\[-2pt]{\tiny $t$ y luego $N$ filas \texttt{x y theta}}};
        \node[archivo,above=2.2mm of dyn] (sta) {\texttt{static.txt}\\[-2pt]{\tiny parámetros de la corrida}};
        \node[archivo,below=2.2mm of dyn] (ob)  {\texttt{observables.txt}\\[-2pt]{\tiny una fila \texttt{t va S} por paso}};

        \node[pyt,right=15mm of dyn] (ani) {\texttt{animate.py}\\[-2pt]{\tiny animaciones y cuadros}};
        \node[pyt,above=2.2mm of ani] (sum) {\texttt{summarise.py}\\[-2pt]{\tiny \teq{} y $\langle X\rangle\pm\sigma_X$}};
        \node[pyt,below=2.2mm of ani] (cur) {\texttt{temporal.py} · \texttt{curves.py}\\[-2pt]{\tiny figuras de los ítems}};

        \begin{scope}[on background layer]
            \node[draw=azul!40,rounded corners=3pt,fit=(arg)(io),inner sep=5pt] (motor) {};
            \node[draw=gris!50,rounded corners=3pt,fit=(sta)(ob),inner sep=5pt] (arch) {};
            \node[draw=acento!45,rounded corners=3pt,fit=(sum)(cur),inner sep=5pt] (post) {};
        \end{scope}

        \node[above=1mm of motor,font=\scriptsize\bfseries,color=azul] {Motor de simulación (C++)};
        \node[above=1mm of arch,font=\scriptsize\bfseries,color=gris] {Archivos de texto};
        \node[above=1mm of post,font=\scriptsize\bfseries,color=acento] {Post-proceso (Python)};

        \draw[fl] (motor.east) -- (arch.west);
        \draw[fl] (arch.east) -- (post.west);
    \end{tikzpicture}
    \end{center}
    \vspace{0.2em}
    {\small El motor nunca dibuja y el post-proceso nunca simula: la velocidad de
    reproducción de la animación no queda atada a la de la simulación.}
\end{frame}

\begin{frame}{El paso de simulación}
    \vspace{-0.1em}
    \begin{columns}[T,onlytextwidth]
        \begin{column}{0.57\textwidth}
            \begin{center}
            \begin{tikzpicture}[
                font=\scriptsize,
                paso/.style={draw=azul!55,fill=azul!6,rounded corners=2pt,align=left,
                             inner sep=4pt,text width=51mm},
                fl/.style={-{Stealth[length=2mm]},azul,thick}]
                \node[paso,fill=grisclaro,draw=gris!60] (e0)
                    {Estado en $t$: \ $\mathbf{r}_i(t)$, \ $\theta_i(t)$};
                \node[paso,below=3mm of e0] (p1)
                    {\textbf{1.} Vecinos $\mathcal{N}_i$ con el CIM};
                \node[paso,below=3mm of p1] (p2)
                    {\textbf{2.} $\Theta_i$ según la regla, más el ruido;\\
                     se escribe en un \alert{buffer aparte}};
                \node[paso,below=3mm of p2] (p3)
                    {\textbf{3.} $\mathbf{r}_i \mathrel{+}= \mathbf{v}_i(t)\,\Delta t$
                     con la dirección \alert{vieja}, y contorno periódico};
                \node[paso,below=3mm of p3,fill=grisclaro,draw=gris!60] (e1)
                    {Se intercambia el buffer: estado en $t+1$};
                \draw[fl] (e0) -- (p1);
                \draw[fl] (p1) -- (p2);
                \draw[fl] (p2) -- (p3);
                \draw[fl] (p3) -- (e1);
                \draw[fl,gris] (e1.east) -- ++(4mm,0) |- ($(e0.east)+(4mm,0)$) -- (e0.east);
            \end{tikzpicture}
            \end{center}
        \end{column}\hfill
        \begin{column}{0.40\textwidth}
            \vspace{1.0em}\small
            \begin{itemize}
                \setlength{\itemsep}{0.9em}
                \item El buffer es lo que hace la actualización \alert{sincrónica}:
                      sin él, las partículas de índice alto se alinearían con
                      direcciones ya actualizadas.
                \item \textbf{CIM}: $\Mc\times\Mc$ celdas; cada partícula se compara
                      sólo contra su celda y las 8 adyacentes. Requiere
                      $L/\Mc > r_c$: con $L = \SI{10}{m}$, $\Mc \le 9$.
                \item Control: la fuerza bruta da trayectorias idénticas con la
                      misma semilla.
            \end{itemize}
        \end{column}
    \end{columns}
\end{frame}

% =========================================================
\section{Simulaciones}
% =========================================================

\begin{frame}{Sistema particular}
    \vspace{-0.4em}
    \begin{columns}[T,onlytextwidth]
        \begin{column}{0.40\textwidth}
            \begin{center}
            \begin{tikzpicture}[scale=0.78,font=\scriptsize,
                part/.style={-{Stealth[length=1.5mm]},thick,gris!85},
                vec/.style={-{Stealth[length=1.7mm]},very thick,azul}]

                \draw[draw=azul!70,thick] (0,0) rectangle (5.6,5.6);
                \foreach \y in {0.9,2.3,3.7,5.1}{
                    \draw[{Stealth[length=1.2mm]}-{Stealth[length=1.2mm]},azul!45]
                        (-0.35,\y) -- (-0.06,\y);
                    \draw[{Stealth[length=1.2mm]}-{Stealth[length=1.2mm]},azul!45]
                        (5.66,\y) -- (5.95,\y);}
                \node[azul!70,anchor=south,font=\scriptsize] at (2.8,5.7) {contorno periódico};
                \draw[|<->|,azul!70] (0,-0.38) -- (5.6,-0.38)
                    node[midway,fill=white,inner sep=1pt] {$L = \SI{10}{m}$};

                \foreach \x/\y/\a in {0.7/4.9/40, 1.6/5.2/10, 4.7/4.9/200,
                    5.1/3.9/170, 0.6/2.4/300, 1.1/0.7/95, 2.4/0.5/120,
                    4.4/0.8/250, 5.2/1.7/25, 3.9/5.1/160, 0.4/3.7/65, 5.2/2.9/210}
                    \draw[part] (\x,\y) -- ++({0.45*cos(\a)},{0.45*sin(\a)});

                \draw[dashed,acento,thick] (2.8,2.9) circle (1.05);
                \draw[acento,{Stealth[length=1.4mm]}-] (2.8,2.9) -- ++(0.74,0.74);
                \node[acento,anchor=south west,inner sep=1pt] at (3.3,3.4) {$r_c = \SI{1}{m}$};
                \foreach \x/\y/\a in {2.2/3.5/58, 3.5/3.2/70, 3.3/2.3/35, 2.3/2.2/80}
                    \draw[vec,azul!70] (\x,\y) -- ++({0.52*cos(\a)},{0.52*sin(\a)});
                \draw[vec,acento] (2.8,2.9) -- ++({0.64*cos(60)},{0.64*sin(60)});
                \node[acento,anchor=north east,inner sep=1pt] at (2.78,2.86) {$i$};
                \node[azul,anchor=west,inner sep=1pt] at (3.95,2.3) {$\mathcal{N}_i$};
            \end{tikzpicture}
            \end{center}
        \end{column}\hfill
        \begin{column}{0.57\textwidth}
            \vspace{0.1em}
            {\scriptsize
            \begin{tabular}{@{}p{3.0cm}l@{}}
                \toprule
                \multicolumn{2}{@{}l}{\color{azul}\bfseries Parámetros fijos}\\
                \midrule
                Lado de la caja $L$        & \SI{10}{m}, contorno periódico \\
                Radio de interacción $r_c$ & \SI{1}{m} \\
                Rapidez $v$                & \SI{0.03}{m/s} \\
                Paso temporal $\Delta t$   & \SI{1}{s} \\
                Celdas del CIM $\Mc$       & $9$ (máximo admisible) \\
                \midrule
                \multicolumn{2}{@{}l}{\color{azul}\bfseries Parámetros variables}\\
                \midrule
                Ruido $\eta$   & $0$ a $\SI{5}{rad}$, paso $\SI{0.25}{rad}$ \\
                Densidad $\rho$ & $2$, $4$, $8\ \si{m^{-2}}$
                                  \ ($N = 200$, $400$, $800$) \\
                                & $\tfrac{1}{\pi}$, $\tfrac{1}{2\pi}$, $\tfrac{1}{3\pi}\ \si{m^{-2}}$
                                  \ ($N = 32$, $16$, $11$) \\
                Regla           & estándar, votante \\
                \midrule
                \multicolumn{2}{@{}l}{\color{azul}\bfseries Corridas}\\
                \midrule
                Realizaciones $M$ & $20$ por caso (semillas 1 a 20) \\
                Duración          & \SI{5e3}{s} si $\rho \ge \SI{2}{m^{-2}}$ \\
                                  & \SI{2e4}{s} si $\rho < \SI{1}{m^{-2}}$ \\
                Total             & \num{5040} corridas \\
                \bottomrule
            \end{tabular}}
        \end{column}
    \end{columns}
\end{frame}

\begin{frame}{Observables}
    \vspace{-0.2em}
    \begin{columns}[T,onlytextwidth]
        \begin{column}{0.53\textwidth}
            {\color{azul}\bfseries Polarización} --- cuán alineado está
            \begin{equation*}
                \va(t) = \frac{1}{N v}\left\|\sum_{i=1}^{N}\mathbf{v}_i(t)\right\|
                \ \in\ [0,1]
            \end{equation*}
            \vspace{-0.7em}

            {\scriptsize Vale 1 si todas van en la misma dirección. Con direcciones
            independientes la suma es una caminata al azar: el piso por tamaño
            finito es $\sqrt{\pi/4N}$, no cero.}
            \vspace{1.0em}

            {\color{azul}\bfseries Componente gigante} --- cuán conectado está
            \begin{equation*}
                S(t) = \frac{|\Cmax(t)|}{N}
                \ \in\ \left[\tfrac{1}{N},\,1\right]
            \end{equation*}
            \vspace{-0.7em}

            {\scriptsize Un nodo por partícula, una arista por cada par a distancia
            menor que $r_c$, y \Cmax la componente conexa más grande.}
        \end{column}\hfill
        \begin{column}{0.43\textwidth}
            \vspace{0.8em}
            \begin{center}
            \begin{tikzpicture}[scale=0.80,font=\scriptsize,
                nodo/.style={circle,fill=gris!45,inner sep=1.5pt},
                gigante/.style={circle,fill=azul,inner sep=1.8pt},
                arista/.style={gris!55,thin},
                aristag/.style={azul!75,semithick}]

                \draw[draw=gris!45] (0,0) rectangle (5.4,4.4);

                \foreach \n/\x/\y in {a/1.0/3.4, b/1.7/3.9, c/1.9/2.9, d/2.7/3.4,
                                      e/2.5/2.4, f/1.2/2.4, g/3.4/2.8}
                    \node[gigante] (\n) at (\x,\y) {};
                \foreach \u/\v in {a/b, a/c, b/c, b/d, c/d, c/e, c/f, e/f, d/e, d/g, e/g}
                    \draw[aristag] (\u) -- (\v);

                \foreach \n/\x/\y in {h/4.4/3.8, i/4.9/3.2, j/0.7/1.0, k/1.4/0.7,
                                      l/1.9/1.2, m/3.6/0.9, n/4.6/1.3}
                    \node[nodo] (\n) at (\x,\y) {};
                \foreach \u/\v in {h/i, j/k, k/l}
                    \draw[arista] (\u) -- (\v);

                \node[azul,anchor=west] at (3.7,3.6) {\Cmax};
                \draw[azul,thin] (3.68,3.55) -- (3.1,3.3);
            \end{tikzpicture}
            \end{center}
            \vspace{0.3em}
            {\scriptsize En el esquema $N = 14$ y $|\Cmax| = 7$, de modo
            que $S = 0.5$. En azul la componente gigante, en gris el resto.}
        \end{column}
    \end{columns}
\end{frame}

\begin{frame}{Del observable temporal al escalar}
    \figpar[0.64]{b_evolucion_temporal/va_S_vs_t_vicsek_rho4_eta2_M20.png}{%
        \setup{Condiciones}
        Modelo estándar\\
        $\rho = \SI{4}{m^{-2}}$ ($N = 400$), $\eta = \SI{2}{rad}$\\
        \fijos\\
        \SI{5e3}{s} por corrida
        \vspace{0.9em}

        \setup{\teq}
        Primer bloque de \SI{50}{s} de la curva \alert{promedio} de $\va$ que cruza
        la media de la segunda mitad. Único por caso, y el mismo para $\va$ y $S$.
        \vspace{0.7em}

        \setup{Escalar y barra}
        Promedio y desvío estándar de \alert{todas} las muestras con $t \ge \teq$
        de las $M$ realizaciones.}
    \vspace{0.3em}
    {\scriptsize
    $\displaystyle
    \langle X\rangle=\frac{1}{MT}\sum_{m=1}^{M}\sum_{t\ge\teq}X_m(t),
    \qquad
    \sigma_X=\sqrt{\frac{1}{MT-1}\sum_{m=1}^{M}\sum_{t\ge\teq}\big(X_m(t)-\langle X\rangle\big)^2},
    \qquad X\in\{\va,\,S\}$}
\end{frame}

% =========================================================
\section{Resultados}
% =========================================================

\begin{frame}{Modelo estándar: animaciones características}
    \dosfig
        {a_animaciones/cuadro_vicsek_rho4_eta0.5_t3000.png}
        {$\eta = \SI{0.5}{rad}$ --- ordenado\\[0.15em]\pendiente}
        {a_animaciones/cuadro_vicsek_rho4_eta4_t3000.png}
        {$\eta = \SI{4}{rad}$ --- desordenado\\[0.15em]\pendiente}
        {\setup{Ambas corridas}
         Modelo estándar\\
         $\rho = \SI{4}{m^{-2}}$ ($N = 400$)\\
         \fijos\\
         semilla 1, cuadro en \SI{3e3}{s}
         \vspace{1.2em}

         Cada flecha es una partícula, con origen en su posición y color según el
         ángulo de su velocidad.
         \vspace{0.7em}

         La longitud es fija: la rapidez es común a todas, así que la flecha indica
         dirección y no módulo.}
\end{frame}

\begin{frame}{Modelo estándar: evolución temporal}
    \figpar[0.72]{b_evolucion_temporal/va_S_vs_t_vicsek_rho4_etas0.5-2-5_semilla1.png}{%
        \setup{Condiciones}
        Modelo estándar\\
        $\rho = \SI{4}{m^{-2}}$ ($N = 400$)\\
        \fijos\\
        \SI{5e3}{s} por corrida\\
        \alert{una realización por curva}
        \vspace{1.0em}

        Tres comportamientos distintos a igual densidad:
        \vspace{0.5em}

        $\eta = \SI{0.5}{rad}$: sube a $\va \simeq 1$ en \SI{200}{s} y casi no fluctúa.
        \vspace{0.4em}

        $\eta = \SI{2}{rad}$: plateau intermedio, fluctuación visible.
        \vspace{0.4em}

        $\eta = \SI{5}{rad}$: nunca ordena.}
\end{frame}

\begin{frame}{Modelo estándar: polarización contra ruido}
    \figpar[0.60]{c_va_vs_eta/va_vs_eta_vicsek_rho2-4-8_M20.png}{%
        \setup{Condiciones}
        Modelo estándar\\
        \fijos\\
        \SI{5e3}{s} por corrida\\
        barra: $\sigma_{\va}$ en el estacionario
        \vspace{1.2em}

        \begin{itemize}\setlength{\itemsep}{0.8em}\setlength{\leftmargini}{1em}
            \item Decae de forma \alert{continua} desde $\va = 1$, sin salto.
            \item A igual $\eta$, \alert{más densidad da más orden}:
                  $\eta(\va = 0.5)$ pasa de $\SI{2.9}{rad}$ a $\SI{3.5}{rad}$
                  entre $\rho = 2$ y $\SI{8}{m^{-2}}$.
            \item Promediar sobre $\langle k\rangle = \rho\,\pi r_c^2$ vecinos
                  atenúa el ruido como $1/\sqrt{\langle k\rangle}$.
        \end{itemize}}
\end{frame}

\begin{frame}{Modelo de votante: animación y evolución temporal}
    \begin{columns}[T,onlytextwidth]
        \begin{column}{0.24\textwidth}
            \centering
            \includegraphics[width=\linewidth]{a_animaciones/cuadro_voter_rho4_eta0.5_t3000.png}\\[0.2em]
            {\scriptsize votante, $\eta = \SI{0.5}{rad}$}\\[0.1em]
            \pendiente
        \end{column}
        \begin{column}{0.50\textwidth}
            \centering
            \includegraphics[width=\linewidth]{f_votante/va_S_vs_t_voter_rho4_etas0.5-2-5_semilla1.png}
        \end{column}\hfill
        \begin{column}{0.23\textwidth}
            \raggedright\scriptsize
            \setup{Condiciones}
            Modelo de votante\\
            $\rho = \SI{4}{m^{-2}}$ ($N = 400$)\\
            \fijos\\
            \SI{5e3}{s} por corrida\\
            una realización por curva
            \vspace{1.2em}

            Con el mismo $\eta$ que ordenaba al estándar, el votante no forma una
            bandada única: aparecen \alert{dominios} de dirección que se arman y se
            disuelven, y $\va$ oscila con amplitud comparable a su media.
        \end{column}
    \end{columns}
\end{frame}

\begin{frame}{Votante contra estándar: polarización contra ruido}
    \dosfiglimpio
        {f_votante/va_vs_eta_vicsek_vs_voter_rho4_M20.png}
        {f_votante/va_vs_eta_voter_rho2-4-8_M20.png}
        {\paneles{ambos modelos, $\rho = \SI{4}{m^{-2}}$}{sólo votante, tres densidades}
         \setup{Ambas figuras}
         \fijos\\
         \SI{5e3}{s} por corrida
         \vspace{0.6em}

         \begin{itemize}\setlength{\itemsep}{0.6em}\setlength{\leftmargini}{1em}
             \item En $\eta = 0$ ambos dan $\va = 1$: verificación cruzada.
             \item El votante cae \alert{un orden de magnitud antes}:
                   $\eta(\va = 0.5) \simeq \SI{0.2}{rad}$ frente a $\SI{3.5}{rad}$.
             \item La densidad tiene \alert{signo opuesto}: copiar a uno no
                   promedia el ruido.
         \end{itemize}}
\end{frame}

\begin{frame}{Clusters: animaciones a baja densidad}
    \dosfig
        {a_animaciones/cuadro_vicsek_rho1pi_eta0.5_t20000.png}
        {$\eta = \SI{0.5}{rad}$ --- agregado, $S \to 1$\\[0.15em]\pendiente}
        {a_animaciones/cuadro_vicsek_rho1pi_eta5_t20000.png}
        {$\eta = \SI{5}{rad}$ --- fragmentado, $S \simeq 0.2$\\[0.15em]\pendiente}
        {\setup{Ambas corridas}
         Modelo estándar\\
         $\rho = 1/\pi\ \si{m^{-2}}$ ($N = 32$)\\
         \fijos\\
         semilla 1, cuadro en \SI{2e4}{s}
         \vspace{1.2em}

         Con $\rho \ge \SI{2}{m^{-2}}$ el grafo de vecinos está percolado siempre y
         $S$ no aporta información.
         \vspace{0.7em}

         A $\rho = 1/\pi$ hay $\langle k\rangle = 1$ vecino por partícula y $S$
         recorre todo su rango.}
\end{frame}

\begin{frame}{Clusters: evolución temporal de la componente gigante}
    \figpar[0.72]{d_clusters/va_S_vs_t_vicsek_rho1pi_etas0.5-2-5_promedioM20.png}{%
        \setup{Condiciones}
        Modelo estándar\\
        $\rho = 1/\pi\ \si{m^{-2}}$ ($N = 32$)\\
        \fijos\\
        \SI{2e4}{s} por corrida\\
        curva promedio de las $M$
        \vspace{1.0em}

        Ahora $S$ tiene \alert{transitorio propio}: parte de $S \simeq 0.15$ y
        crece por agregación.
        \vspace{0.7em}

        Dos partículas que se encuentran se alinean, y al alinearse dejan de
        separarse: la unión sólo se deshace por el ruido.
        \vspace{0.7em}

        \teq{} llega a \SI{3.65e3}{s}.}
\end{frame}

\begin{frame}{Clusters: componente gigante contra ruido}
    \dosfiglimpio
        {d_clusters/S_vs_eta_vicsek_rho8-2-1pi-1-3pi_M20.png}
        {d_clusters/S_vs_eta_voter_rho8-2-1pi-1-3pi_M20.png}
        {\paneles{modelo estándar}{modelo de votante}
         \setup{Ambas figuras}
         \fijos\\
         \SI{5e3}{s} si $\rho \ge \SI{2}{m^{-2}}$, \SI{2e4}{s} si no\\
         mismo eje $[0,1]$ en las dos
         \vspace{0.4em}

         \begin{itemize}\setlength{\itemsep}{0.35em}\setlength{\leftmargini}{1em}
             \item $\rho = 8$ y $2$ quedan \alert{pegadas a $S = 1$}:
                   $\langle k\rangle = 25$ y $6{,}3$, muy por encima del umbral
                   ($\simeq 4{,}5$).
             \item $1/\pi$ y $1/3\pi$ \alert{decaen monótonamente} hasta
                   $S \simeq 0.2$.
             \item En $\eta = 0$ percolan igual: \alert{es el ruido lo que
                   fragmenta}.
         \end{itemize}}
\end{frame}

\begin{frame}{Polarización contra componente gigante}
    \dosfiglimpio
        {e_va_vs_S/va_vs_S_vicsek_rho8-2-1pi-1-3pi_M20.png}
        {e_va_vs_S/va_vs_S_voter_rho8-2-1pi-1-3pi_M20.png}
        {\paneles{modelo estándar}{modelo de votante}
         \setup{Ambas figuras}
         Un punto por cada $\eta$ del barrido\\
         \fijos\\
         barras: $\sigma_S$ y $\sigma_{\va}$
         \vspace{0.4em}

         \begin{itemize}\setlength{\itemsep}{0.35em}\setlength{\leftmargini}{1em}
             \item A $\rho = 8$ y $2$ los puntos se apilan contra $S = 1$:
                   \alert{alineación y conectividad, desacopladas}.
             \item A $1/\pi$ y $1/3\pi$ la relación es \alert{monótona creciente};
                   las dos casi colapsan.
             \item Pearson: $0{,}45$ y $0{,}61$ frente a $0{,}985$ y $0{,}994$.
         \end{itemize}}
\end{frame}

\begin{frame}{Tiempos de ejecución del CIM}
    \figpar[0.60]{g_tiempos_cim/tiempo_busqueda_vs_N_L20_M13_cim_bruta_tp1_log.png}{%
        \setup{Caja del TP1}
        $L = \SI{20}{m}$, $r_c = \SI{1}{m}$\\
        $\Mc = 13$, \SI{200}{s} por corrida\\
        5 repeticiones por punto\\
        barra: desvío entre repeticiones\\
        equipo descargado
        \vspace{1.2em}

        \begin{itemize}\setlength{\itemsep}{0.8em}\setlength{\leftmargini}{1em}
            \item Exponentes ajustados: $2{,}11$ (fuerza bruta) y $1{,}34$ (CIM).
                  No es 1 porque $L$ es fijo: al subir $N$ sube $\rho$.
            \item Lo que baja el CIM es el \alert{prefactor}: cada partícula se
                  compara contra el $9/\Mc^2 = \SI{5.3}{\percent}$ del sistema.
            \item Cruce en $N \simeq 90$; en $N = 1150$ el CIM es \num{10} veces
                  más rápido.
            \item Misma pendiente que el TP1; el corrimiento vertical es del equipo.
        \end{itemize}}
\end{frame}

% =========================================================
\section{Conclusiones}
% =========================================================

\begin{frame}{Conclusiones}
    \vspace{-0.2em}
    \begin{columns}[T,onlytextwidth]
        \begin{column}{0.48\textwidth}\small
            \begin{itemize}\setlength{\itemsep}{0.9em}
                \item El modelo estándar exhibe una \alert{transición continua}
                      orden--desorden controlada por $\eta$, y \alert{la densidad
                      favorece el orden}: $\eta(\va = 0.5)$ pasa de $\SI{2.9}{rad}$
                      a $\SI{3.5}{rad}$ entre $\rho = 2$ y $\SI{8}{m^{-2}}$.
                \item El votante es \alert{un orden de magnitud más sensible} al
                      ruido y la densidad le juega con \alert{signo opuesto}. En
                      $\eta = 0$ ambos coinciden en $\va = 1$.
                \item En $\eta = \SI{5}{rad}$ el votante alcanza el piso
                      $\sqrt{\pi/4N}$ en las cuatro densidades; el estándar no.
            \end{itemize}
        \end{column}\hfill
        \begin{column}{0.48\textwidth}\small
            \begin{itemize}\setlength{\itemsep}{0.9em}
                \item Con $\rho = 2$, $4$, $\SI{8}{m^{-2}}$ la componente gigante
                      \alert{no es un grado de libertad}: $S \ge 0.90$ siempre y
                      $\va$ no es función de $S$.
                \item Al bajar a $\rho = 1/\pi$ y $1/3\pi$ ese desacople desaparece
                      y $S$ queda \alert{fuertemente acoplado} a $\va$: es un efecto
                      de los parámetros, no del modelo.
                \item Sin ruido el sistema percola aun con un vecino por partícula:
                      \alert{es el ruido lo que fragmenta}.
                \item El CIM es \num{10} veces más rápido que la fuerza bruta en
                      $N = 1150$, con el mismo escalamiento que en el TP1.
            \end{itemize}
        \end{column}
    \end{columns}
\end{frame}

\begin{frame}[plain,noframenumbering]
    \vfill
    \begin{center}
        {\Large\bfseries\color{azul} Muchas gracias}\\[0.6em]
        \textcolor{azul!30}{\rule{0.25\textwidth}{1pt}}
    \end{center}
    \vfill
\end{frame}

\end{document}
